\section*{Exact GU accounting for antichain addressing with arbitrary reuse}

All logarithms are base two.  Let $Y=(y_1,\ldots,y_k)$ be selector
coordinates, let $C$ be a set of common data coordinates, and let
$\{A_z:z\in\{-1,1\}^k\}$ be an antichain of distinct subsets of the
remaining physical data coordinates.  No bound is imposed on the sizes,
pairwise intersections, or congestion of the $A_z$.  Define
\[
 F(Y,x)=\chi_C(x)\chi_{A_Y}(x).                       \tag{1}
\]
This is a genuine Boolean function and includes full physical reuse.

\paragraph{Theorem.}
Put $m=2^k$ and, for a data coordinate $i$, set
\[
 B_{i,z}=A_z\setminus\{i\},\qquad
 d_i(z)=\#\{w:B_{i,z}\subseteq B_{i,w}\}.            \tag{2}
\]
Then
\[
 U(F)=\frac{k}{2},\qquad
 I(F)=|C|+\frac{k}{2}+\frac1m\sum_z|A_z|,             \tag{3}
\]
and the complete signed GU aggregate is
\[
 \boxed{G(F)=\frac{k}{2}-\frac1m\sum_{i,z}\log_2 d_i(z).} \tag{4}
\]
Consequently this entire all-arity, unbounded-reuse class obeys
\[
 \boxed{G(F)\le I(F)}                                 \tag{5}
\]
with the dimension-independent restricted GU pair $(a,b)=(1,0)$.

\paragraph{Proof.}
Writing $1_{\{Y=z\}}=m^{-1}\sum_{R\subseteq[k]}\sigma_z(R)\chi_R(Y)$
shows that (1) has $m^2$ Fourier coefficients of squared mass $m^{-2}$,
at supports $C\cup A_z\cup R$.  The antichain assumption makes all these
supports distinct and prevents cross-branch upper relations.  For fixed
$(z,R)$ there are $2^{k-|R|}$ upper supersets, hence
\[
 \log_2\frac{\pi(C\cup A_z\cup R)}{p(C\cup A_z\cup R)}=k-|R|.
\]
Averaging over uniform $R$ proves $U=k/2$.  Averaging support cardinality
proves the formula for $I$.

Restricting a selector $y_j$ leaves an antichain address on $k-1$
selectors in either child.  Thus every selector contributes exactly
$\Delta_{y_j}U=1/2$.  Restricting a coordinate in $C$ only removes a
common character and contributes zero.

For a remaining data coordinate $i$, antichainness also implies that the
$B_{i,z}$ are distinct: equality would make two original $A_z$ equal or
nested.  The restriction signs therefore cannot merge Fourier
coefficients, so both children have the same squared spectrum.  A
coefficient indexed by $(z,R)$ now sees exactly $d_i(z)$ eligible branch
supports and $2^{k-|R|}$ selector supersets.  Hence each child has
\[
 U(F_{i,\pm})=\frac{k}{2}+\frac1m\sum_z\log_2d_i(z),
\]
and therefore
\[
 \Delta_iU=-\frac1m\sum_z\log_2d_i(z).                \tag{6}
\]
Summing the $k$ selector increments, the common-coordinate zeros, and
(6) proves (4).  Since every $d_i(z)\ge1$, (4) gives $G\le k/2\le I$.

\paragraph{Interpretation and scope.}
Equation (6) is the complete cost of coherent upper-set overlap created
by reusing a physical input: every new containment is a signed refund,
with its exact logarithmic multiplicity retained.  No disjointness,
bounded congestion, cancellation assumption, or correlation gap is
used.  The theorem extends the previous two-piece collision calculation
to arbitrary addressing depth and arbitrary antichain overlap.

It does not prove universal GU.  The remaining obstruction is the first
stage at which the branch Fourier rows are not single characters (or
where initial branch supports are comparable): Fourier coefficients can
then merge after restriction, so the signed amplitudes, rather than only
the containment counts $d_i(z)$, determine the posterior cost.
